\documentclass[10pt,a4paper]{amsart}
\usepackage{amsmath,amssymb,amsthm,graphicx}
\usepackage{verbatim}
\newenvironment{reference}{\comment}{\endcomment}
%\newenvironment{reference}{}{}
\newenvironment{slogan}{\comment}{\endcomment}
\newenvironment{history}{\comment}{\endcomment}
% For commutative diagrams we use Xy-pic
\usepackage[all]{xy}

% We use 2cell for 2-commutative diagrams.
\xyoption{2cell}
% Theorem environments.
%
\theoremstyle{plain}
\newtheorem{theorem}[subsection]{Theorem}
\newtheorem{proposition}[subsection]{Proposition}
\newtheorem{lemma}[subsection]{Lemma}

\theoremstyle{definition}
\newtheorem{definition}[subsection]{Definition}
\newtheorem{example}[subsection]{Example}
\newtheorem{exercise}[subsection]{Exercise}
\newtheorem{situation}[subsection]{Situation}

\theoremstyle{remark}
\newtheorem{remark}[subsection]{Remark}
\newtheorem{remarks}[subsection]{Remarks}

\numberwithin{equation}{subsection}

% Macros
%
\def\lim{\mathop{\mathrm{lim}}\nolimits}
\def\colim{\mathop{\mathrm{colim}}\nolimits}
\def\Spec{\mathop{\mathrm{Spec}}}
\def\Hom{\mathop{\mathrm{Hom}}\nolimits}
\def\Ext{\mathop{\mathrm{Ext}}\nolimits}
\def\SheafHom{\mathop{\mathcal{H}\!\mathit{om}}\nolimits}
\def\SheafExt{\mathop{\mathcal{E}\!\mathit{xt}}\nolimits}
\def\Sch{\mathit{Sch}}
\def\Mor{\mathop{\mathrm{Mor}}\nolimits}
\def\Ob{\mathop{\mathrm{Ob}}\nolimits}
\def\Sh{\mathop{\mathit{Sh}}\nolimits}
\def\NL{\mathop{N\!L}\nolimits}
\def\CH{\mathop{\mathrm{CH}}\nolimits}
\def\proetale{{pro\text{-}\acute{e}tale}}
\def\etale{{\acute{e}tale}}
\def\QCoh{\mathit{QCoh}}
\def\Ker{\mathop{\mathrm{Ker}}}
\def\Im{\mathop{\mathrm{Im}}}
\def\Coker{\mathop{\mathrm{Coker}}}
\def\Coim{\mathop{\mathrm{Coim}}}

% Boxtimes
%
\DeclareMathSymbol{\boxtimes}{\mathbin}{AMSa}{"02}

%
% Macros for moduli stacks/spaces
%
\def\QCohstack{\mathcal{QC}\!\mathit{oh}}
\def\Cohstack{\mathcal{C}\!\mathit{oh}}
\def\Spacesstack{\mathcal{S}\!\mathit{paces}}
\def\Quotfunctor{\mathrm{Quot}}
\def\Hilbfunctor{\mathrm{Hilb}}
\def\Curvesstack{\mathcal{C}\!\mathit{urves}}
\def\Polarizedstack{\mathcal{P}\!\mathit{olarized}}
\def\Complexesstack{\mathcal{C}\!\mathit{omplexes}}
% \Pic is the operator that assigns to X its picard group, usage \Pic(X)
% \Picardstack_{X/B} denotes the Picard stack of X over B
% \Picardfunctor_{X/B} denotes the Picard functor of X over B
\def\Pic{\mathop{\mathrm{Pic}}\nolimits}
\def\Picardstack{\mathcal{P}\!\mathit{ic}}
\def\Picardfunctor{\mathrm{Pic}}
\def\Deformationcategory{\mathcal{D}\!\mathit{ef}}
\usepackage{cleveref}
\crefname{originallemma}{original lemma}{original lemmas}
\Crefname{originallemma}{Original lemma}{Original lemmas}
\makeatletter
\newlabel{algebra-lemma-finite-in-codim-1}{{02MA}{source}{}{sourcepointer-1}{}}
\newlabel{algebra-lemma-finite-in-codim-1@cref}{{[originallemma][0][]02MA}{[1][0][]source}{}{}{}}
\makeatother
\crefname{originallemma}{original lemma}{original lemmas}
\Crefname{originallemma}{Original lemma}{Original lemmas}
\makeatletter
\newlabel{coherent-lemma-proper-finite-fibre-finite-in-neighbourhood}{{02OH}{source}{}{sourcepointer-2}{}}
\newlabel{coherent-lemma-proper-finite-fibre-finite-in-neighbourhood@cref}{{[originallemma][0][]02OH}{[1][0][]source}{}{}{}}
\makeatother
\crefname{originallemma}{original lemma}{original lemmas}
\Crefname{originallemma}{Original lemma}{Original lemmas}
\makeatletter
\newlabel{varieties-lemma-finite-in-codim-1}{{0AB7}{source}{}{sourcepointer-3}{}}
\newlabel{varieties-lemma-finite-in-codim-1@cref}{{[originallemma][0][]0AB7}{[1][0][]source}{}{}{}}
\makeatother
\title{A proper dominant morphism with finite function field extension is finite over a neighbourhood of a codimension one point}
\author{The Stacks Project Authors}
\date{stacks-project:chow.tex:local-snapshot-20261012; GNU FDL 1.2 or later}
\begin{document}
\maketitle
\noindent\textit{Editorial note.} All delivered mathematical text is the author's own, copied from the pinned Stacks Project chapter file chow.tex; no statement or proof step was rewritten, reordered or removed.\par\medskip
\noindent\textit{Editorial note.} Preamble provenance: the pinned chapter file opens with \textbackslash{}input\{preamble\} and carries no preamble text of its own, so the project's single Stacks-wide preamble staged at sources/stacks-chow/preamble.tex is excerpted verbatim as three spans, lines 8-12 (the verbatim package together with the project's own reference, slogan and history display-environment definitions, including the author's commented-out older form of the reference environment), lines 14-18 (Xy-pic setup) and lines 154-217 (the native theorem declarations and mathematical macros). The lines 8-12 span is required because a counted statement or its surrounding text may use those display environments, which the project typesets as comments. Nothing was added, removed or redefined; the wrapper class and the amsmath/amssymb/amsthm/graphicx base packages are supplied by the assembler.\par\medskip
\noindent\textit{Editorial note.} Same-chapter pointers are located at their \textbackslash{}label lines inside the pinned chapter file; cross-chapter pointers are located in the project's own tag list staged in the snapshot (sources/stacks-chow/tags.txt), at the verbatim lines carrying their canonical tags.\par\medskip
\noindent\textit{Allowed standard prerequisites (editorial).}\begin{itemize}
\item Properness with finite fibres in a neighbourhood, and the algebraic criterion in codimension one. Cohomology of Schemes, label coherent-lemma-proper-finite-fibre-finite-in-neighbourhood (proof line 2543); Algebra, label algebra-lemma-finite-in-codim-1 (proof line 2553). The first reduces the statement to the finiteness of the fibre over the point, and the second provides the algebraic criterion giving finiteness over a neighbourhood in the one-dimensional case.
\end{itemize}\medskip
\ifdefined\hypertarget\hypertarget{sourcepointer-1}{}\fi
\noindent\textit{Source reference.} Original lemma 02MA: sources/stacks-chow/tags.txt, source lines 3437--3437. Navigation pointer for \textbackslash{}ref\{algebra-lemma-finite-in-codim-1\} used inside the counted span (line 2553); the argument that consumes it is retained in the span.\par\medskip
\ifdefined\hypertarget\hypertarget{sourcepointer-2}{}\fi
\noindent\textit{Source reference.} Original lemma 02OH: sources/stacks-chow/tags.txt, source lines 3516--3516. Navigation pointer for \textbackslash{}ref\{coherent-lemma-proper-finite-fibre-finite-in-neighbourhood\} used inside the counted span (line 2543); the argument that consumes it is retained in the span.\par\medskip
\ifdefined\hypertarget\hypertarget{sourcepointer-3}{}\fi
\noindent\textit{Source reference.} Original lemma 0AB7: sources/stacks-chow/tags.txt, source lines 13010--13010. Navigation pointer for \textbackslash{}ref\{varieties-lemma-finite-in-codim-1\} used inside the counted span (line 2540); the argument that consumes it is retained in the span.\par\medskip
\begin{lemma}
\label{lemma-finite-in-codimension-one}
Let $f : X \to Y$ be a morphism of schemes.
Let $\xi \in Y$ be a point.
Assume that
\begin{enumerate}
\item $X$, $Y$ are integral,
\item $Y$ is locally Noetherian
\item $f$ is proper, dominant and $R(Y) \subset R(X)$ is finite, and
\item $\dim(\mathcal{O}_{Y, \xi}) = 1$.
\end{enumerate}
Then there exists an open neighbourhood $V \subset Y$ of $\xi$
such that $f|_{f^{-1}(V)} : f^{-1}(V) \to V$ is finite.
\end{lemma}

\begin{proof}
This lemma is a special case of
Varieties, Lemma \ref{varieties-lemma-finite-in-codim-1}.
Here is a direct argument in this case.
By Cohomology of Schemes,
Lemma \ref{coherent-lemma-proper-finite-fibre-finite-in-neighbourhood}
it suffices to prove that $f^{-1}(\{\xi\})$ is finite.
We replace $Y$ by an affine open, say $Y = \Spec(R)$.
Note that $R$ is Noetherian, as $Y$ is assumed locally Noetherian.
Since $f$ is proper it is quasi-compact. Hence we can find a finite
affine open covering $X = U_1 \cup \ldots \cup U_n$ with
each $U_i = \Spec(A_i)$. Note that $R \to A_i$ is a
finite type injective homomorphism of domains such that
the induced extension of fraction fields is finite.
Thus the lemma follows
from Algebra, Lemma \ref{algebra-lemma-finite-in-codim-1}.
\end{proof}

\end{document}
